\scanpage{87}
\mspage{11}
alors que pratiquement jamais on ne pourra trouver un $x\in X_K$ \deleted{stable} \inserted{fixe} sous $\unsure{\mathcal G|G}$.
\figure{\begin{tikzpicture}\draw (0,0)--(1.4,0);\end{tikzpicture}}

Avant de donner un principe d'approche pour des calculs ``symétriques'', tenant compte de la présence d'un groupe $\unsure{\mathcal G|G}$ (cas d'une grande importance en pratique), je voudrais préciser la question des relations entre $\Pi_1(X,I)$ ($I$ partie finie de $X_K$) et $\Pi_1(U_{0,3})$, \deleted{et les} dans l'optique des opérations de $\Gamma$.
\inserted{Pour fixer les idées je suppose que $K$ est \deleted{de} fini sur $\mathbb Q$,} \inserted{et} Je prendrai le cas $X$ quasi-projective lisse géom.$^t$ connexe. Utilisant des th. du type Lefschetz, plus le résultat de \deleted{\unsure{Biel}} \unsure{Белый|Belyi}, on \deleted{\unsure{trouve|trouver}} doit pouvoir trouver deux hom de schémas sur $K$
\[
\tag{44}\diagram{\begin{tikzcd}[ampersand replacement=\&]
I \arrow[d,phantom,"{\rotatebox{-90}{$\subset$}}" description] \& \\
\unsure{\mathcal C|C} \arrow[r,hook,"i"] \arrow[d,"f"'] \& X \\
U_{0,3} \&
\end{tikzcd}}
\]
où $i$ induit un hom \uline{surjectif} sur les groupes fondamentaux, et où $f$ est
\margin{86}
